\numoptsfalse
\yearpaper{2024}{Day 2 \textperiodcentered{} First Session \textperiodcentered{} Version A-2 \textperiodcentered{} 60 Questions \textperiodcentered{} 60 Marks \textperiodcentered{} 80 Minutes}{}
\begin{multicols}{2}

\QQ{An equiconvex lens of radius of curvature $14\unit{cm}$ is made up of two different materials. Left half and right half of vertical portion is made up of material of refractive index 1.5 and 1.2 respectively as shown in the figure. If a point object is placed at a distance of $40\unit{cm}$, calculate the image distance.
\fig{\begin{tikzpicture}[scale=0.8]
\fill[pattern=north east lines] (0,1.5) .. controls (-0.55,0.5) and (-0.55,-0.5) .. (0,-1.5) -- cycle;
\draw (0,1.5) .. controls (0.55,0.5) and (0.55,-0.5) .. (0,-1.5) .. controls (-0.55,-0.5) and (-0.55,0.5) .. (0,1.5);
\draw[dashed] (0,1.5) -- (0,-1.5);
\draw[<-] (-0.25,0.3) -- (-0.9,0.3) node[left,lbl]{1.5};
\draw[<-] (0.25,0.3) -- (0.9,0.3) node[right,lbl]{1.2};
\end{tikzpicture}}}
{\opt{$25\unit{cm}$}{$50\unit{cm}$}{$35\unit{cm}$}{$40\unit{cm}$}}

\QQ{A galaxy is moving away from the Earth so that a spectral line at $600\unit{nm}$ is observed at $601\unit{nm}$. Then the speed of the galaxy with respect to the Earth is}
{\opt{$500\unit{km\,s^{-1}}$}{$50\unit{km\,s^{-1}}$}{$200\unit{km\,s^{-1}}$}{$20\unit{km\,s^{-1}}$}}

\QQ{Three polaroid sheets are co-axially placed as indicated in the diagram. Pass axes of the polaroids 2 and 3 make $30^\circ$ and $90^\circ$ with pass axis of polaroid sheet 1. If $I_0$ is the intensity of the incident unpolarised light entering sheet 1, the intensity of the emergent light through sheet 3 is
\fig{\begin{tikzpicture}[scale=0.62]
\draw[dashed] (-1.6,0) -- (5.6,0);
\draw[->] (-1.4,0.35) -- (-0.6,0.35) node[above,pos=0.4,lbl]{$I_0$};
\foreach \x/\n in {0/1,2/2,4/3}{\draw (\x,0) ellipse (0.45 and 1.1); \node[lbl] at (\x+0.35,-1.2) {\n};}
\draw[densely dashed] (0,-1.5) -- (0,1.6) node[above right,lbl,inner sep=1pt]{PA$_1$};
\draw[densely dashed] (2,0) -- ($(2,0)+(60:1.6)$) node[above right,lbl,inner sep=1pt]{PA$_2$}; \draw[densely dashed] (2,0) -- ($(2,0)+(240:1.5)$);
\draw (2,1.25) arc (90:62:1.25); \node[lbl] at (2.35,1.55) {$30^\circ$};
\draw[densely dashed] (4,0) -- ($(4,0)+(20:1.6)$) node[right,lbl,inner sep=1pt]{PA$_3$}; \draw[densely dashed] (4,0) -- ($(4,0)+(200:1.5)$);
\draw[densely dashed] (4,-1.5) -- (4,1.6);
\node[lbl] at (4.6,1.1) {$90^\circ$};
\end{tikzpicture}}}
{\opt{zero}{$\dfrac{3I_0}{32}$}{$\dfrac{3I_0}{8}$}{$\dfrac{3I_0}{16}$}}

\QQ{In Young's double slit experiment, an electron beam is used to produce interference fringes of width $\beta_1$. Now the electron beam is replaced by a beam of protons with the same experimental set-up and same speed. The fringe width obtained is $\beta_2$. The correct relation between $\beta_1$ and $\beta_2$ is}
{\opt{$\beta_1 = \beta_2$}{No fringes are formed}{$\beta_1 < \beta_2$}{$\beta_1 > \beta_2$}}

\QQ{Light of energy $E$ falls normally on a metal of work function $\dfrac E3$. The kinetic energies ($K$) of the photo electrons are}
{\opt{$K = \dfrac{2E}{3}$}{$K = \dfrac{E}{3}$}{$0\le K\le\dfrac{2E}{3}$}{$0\le K\le\dfrac{E}{3}$}}

\QQ{The photoelectric work function for photo metal is $2.4\unit{eV}$. Among the four wavelengths, the wavelength of light for which photo-emission does not take place is}
{\opt{$200\unit{nm}$}{$300\unit{nm}$}{$700\unit{nm}$}{$400\unit{nm}$}}

\QQ{In alpha particle scattering experiment, if $v$ is the initial velocity of the particle, then the distance of closest approach is $d$. If the velocity is doubled, then the distance of closest approach becomes}
{\opt{$4d$}{$2d$}{$\dfrac d2$}{$\dfrac d4$}}

\QQ{The ratio of area of first excited state to ground state of orbit of hydrogen atom is}
{\opt{1:16}{1:4}{4:1}{16:1}}

\QQ{The ratio of volume of $\mathrm{Al}^{27}$ nucleus to its surface area is (Given $R_0 = 1.2\times10^{-15}\unit{m}$)}
{\opt{$2.1\times10^{-15}\unit{m}$}{$1.3\times10^{-15}\unit{m}$}{$0.22\times10^{-15}\unit{m}$}{$1.2\times10^{-15}\unit{m}$}}

\QQ{Consider the nuclear fission reaction
\[{}^1_0\mathrm n + {}^{235}_{92}\mathrm U \to {}^{144}_{56}\mathrm{Ba} + {}^{89}_{36}\mathrm{Kr} + 3\,{}^1_0\mathrm n.\]
Assuming all the kinetic energy is carried away by the fast neutrons only and total binding energies of ${}^{235}_{92}\mathrm U$, ${}^{144}_{56}\mathrm{Ba}$ and ${}^{89}_{36}\mathrm{Kr}$ to be $1800\unit{MeV}$, $1200\unit{MeV}$ and $780\unit{MeV}$ respectively, the average kinetic energy carried by each fast neutron is (in MeV)}
{\opt{200}{180}{67}{60}}

\QQ{The natural logarithm of the activity $R$ of a radioactive sample varies with time $t$ as shown. At $t=0$, there are $N_0$ undecayed nuclei. Then $N_0$ is equal to [Take $e^2 = 7.5$]
\fig{\begin{tikzpicture}[scale=0.8]
\draw[->] (-0.2,0) -- (3.6,0) node[right,lbl]{$t$ (in $10^3$ s)};
\draw[->] (0,-0.2) -- (0,2.6) node[above,lbl]{$\ln R$};
\draw (0,2) -- (2.4,0.25);
\fill (0,2) circle (1.5pt) (0,1) circle (1.5pt);
\node[lbl,left] at (0,2) {2}; \node[lbl,left] at (0,1) {1};
\draw[dashed] (0,1) -- (1.37,1) -- (1.37,0) node[below,lbl]{10};
\end{tikzpicture}}}
{\opt{7,500}{3,500}{75,000}{1,50,000}}

\QQ{Depletion region in an unbiased semiconductor diode is a region consisting of}
{\opt{both free electrons and holes}{neither free electrons nor holes}{only free electrons}{only holes}}

\QQ{The upper level of valence band and lower level of conduction band overlap in the case of}
{\opt{silicon}{copper}{carbon}{germanium}}

\QQ{In the diagram shown, the Zener diode has a reverse breakdown voltage of $V_Z$. The current through the load resistance $R_L$ is $I_L$. The current through the Zener diode is
\fig{\begin{circuitikz}[scale=0.62,transform shape]
\draw (0,0) to[battery1,l=$V_0$,invert] (0,3.4) -- (2,3.4) to[R,l=$R_S$] (2,1.9);
\draw (2,0) to[zDo,i<^=$I_Z$] (2,1.9) -- (4,1.9) to[R,l=$R_L$,i>^=$I_L$] (4,0) -- (0,0);
\end{circuitikz}}}
{\opt{$\dfrac{V_0-V_Z}{R_S}$}{$\dfrac{V_0-V_Z}{R_L}$}{$\dfrac{V_Z}{R_L}$}{$\left(\dfrac{V_0-V_Z}{R_S}\right)-I_L$}}

\QQ{A p-n junction diode is connected to a battery of emf $5.7\unit{V}$ in series with a resistance $5\unit{k\Omega}$ such that it is forward biased. If the barrier potential of the diode is $0.7\unit{V}$, neglecting the diode resistance, the current in the circuit is}
{\opt{$1.14\unit{mA}$}{$1\unit{mA}$}{$1\unit{A}$}{$1.14\unit{A}$}}

\QQ{A block of certain mass is placed on a rough inclined plane. The angle between the plane and the horizontal is $30^\circ$. The coefficients of static and kinetic frictions between the block and the inclined plane are 0.6 and 0.5 respectively. Then the magnitude of the acceleration of the block is [Take $g = 10\unit{m\,s^{-2}}$]
\fig{\begin{tikzpicture}[scale=0.8]
\fill[pattern=north east lines] (0,-0.12) rectangle (3.2,0);
\draw (0,0) -- (3.2,0) -- (3.2,1.85) -- cycle;
\begin{scope}[shift={(1.3,0.75)},rotate=30]
\draw[pattern=north east lines] (0,0) rectangle (0.7,0.5);
\end{scope}
\draw (0.6,0) arc (0:30:0.6); \node[lbl] at (0.85,0.2) {$30^\circ$};
\end{tikzpicture}}}
{\opt{$2\unit{m\,s^{-2}}$}{zero}{$0.196\unit{m\,s^{-2}}$}{$0.67\unit{m\,s^{-2}}$}}

\QQ{A particle of mass $500\unit{g}$ is at rest. It is free to move along a straight line. The power delivered to the particle varies with time according to the following graph:
\fig{\begin{tikzpicture}[xscale=0.5,yscale=0.18]
\draw[->] (0,0) -- (6.6,0) node[right,lbl]{$t$ (s)};
\draw[->] (0,0) -- (0,13.5) node[above,lbl]{$P$ (watt)};
\foreach \x in {1,...,6} \node[lbl,below] at (\x,0) {\x};
\foreach \y in {2,4,...,12} {\draw (0.1,\y) -- (-0.1,\y) node[left,lbl]{\y};}
\draw (0,2) -- (6,9.2);
\draw[dashed] (0,8) -- (5,8) -- (5,0);
\fill (5,8) circle (4pt);
\end{tikzpicture}}
The momentum of the particle at $t = 5\unit{s}$ is}
{\opt{$2\sqrt5\unit{N\,s}$}{$5\sqrt2\unit{N\,s}$}{$5\unit{N\,s}$}{$5.5\unit{N\,s}$}}

\QQ{Dimensional formula for activity of a radioactive substance is}
{\opt{$[\mathrm{M^0L^1T^{-1}}]$}{$[\mathrm{M^0L^{-1}T^0}]$}{$[\mathrm{M^0L^0T^{-1}}]$}{$[\mathrm{M^{-1}L^0T^0}]$}}

\QQ{An athlete runs along a circular track of diameter $80\unit{m}$. The distance travelled and the magnitude of displacement of the athlete when he covers $\frac34$th of the circle is (in m)}
{\opt{$60\pi$, $40\sqrt2$}{$40\pi$, $60\sqrt2$}{$120\pi$, $80\sqrt2$}{$80\pi$, $120\sqrt2$}}

\QQ{Among the given pair of vectors, the resultant of two vectors can never be 3 units. The vectors are}
{\opt{1 unit and 2 units}{2 units and 5 units}{3 units and 6 units}{4 units and 8 units}}

\QQ{A ceiling fan is rotating around a fixed axle as shown. The direction of angular velocity is along \rule{1.2cm}{0.4pt}.
\fig{\begin{tikzpicture}[scale=0.7]
\draw[densely dashed] (0,-1.6) -- (0,2.4);
\draw (-0.2,0.3) -- (-0.2,2.0) (0.2,0.3) -- (0.2,2.0); \draw (-0.2,2.0) -- (0.2,2.0);
\draw (-0.2,0.3) -- (-0.55,-0.15) -- (0.55,-0.15) -- (0.2,0.3);
\draw (-0.55,-0.15) -- (-0.55,-0.4) arc (180:360:0.55 and 0.15) -- (0.55,-0.15);
\draw[rotate around={-40:(-0.4,-0.3)}] (-0.4,-0.45) rectangle (-2.1,-0.15);
\draw[rotate around={40:(0.4,-0.3)}] (0.4,-0.45) rectangle (2.1,-0.15);
\draw[rotate around={-150:(-0.3,0.0)}] (-0.3,-0.12) rectangle (-1.9,0.18);
\draw[->] (0.28,-1.05) arc (0:-300:0.28 and 0.12);
\draw[densely dashed,->] (2.2,1.0) -- (2.2,2.0) node[right,lbl]{Z};
\draw[densely dashed,->] (2.2,1.0) -- (3.3,1.0) node[right,lbl]{Y};
\draw[densely dashed,->] (2.2,1.0) -- (1.6,0.4) node[below,lbl]{X};
\end{tikzpicture}}}
{\opt{$+\hat j$}{$-\hat j$}{$+\hat k$}{$-\hat k$}}

\QQ{A body of mass $1\unit{kg}$ is suspended by a weightless string which passes over a frictionless pulley of mass $2\unit{kg}$ as shown in the figure. The mass is released from a height of $1.6\unit{m}$ from the ground. With what velocity does it strike the ground?
\fig{\begin{tikzpicture}[scale=0.75]
\draw (0,2.6) circle (0.55); \fill (0,2.6) circle (1.3pt); \node[lbl] at (0,2.85) {2 kg};
\draw (0,2.6) -- (1.4,2.6); \fill[pattern=north east lines] (1.4,1.9) rectangle (1.52,3.2); \draw (1.4,1.9) -- (1.4,3.2);
\draw (-0.55,2.6) -- (-0.55,0.9);
\draw[pattern=north east lines] (-0.75,0.55) rectangle (-0.35,0.9); \node[lbl,right] at (-0.3,0.72) {1 kg};
\draw[<->] (-0.55,0.55) -- (-0.55,0) node[midway,right,lbl]{1.6 m};
\fill[pattern=north east lines] (-1.4,-0.12) rectangle (1.6,0); \draw (-1.4,0) -- (1.6,0);
\end{tikzpicture}}}
{\opt{$16\unit{m\,s^{-1}}$}{$8\unit{m\,s^{-1}}$}{$4\sqrt2\unit{m\,s^{-1}}$}{$4\unit{m\,s^{-1}}$}}

\QQ{What is the value of acceleration due to gravity at a height equal to half the radius of the Earth, from its surface?}
{\opt{$4.4\unit{m\,s^{-2}}$}{$6.5\unit{m\,s^{-2}}$}{zero}{$9.8\unit{m\,s^{-2}}$}}

\QQ{A thick metal wire of density $\rho$ and length `$L$' is hung from a rigid support. The increase in length of the wire due to its own weight is ($Y$ = Young's modulus of the material of the wire)}
{\opt{$\dfrac{\rho gL}{Y}$}{$\dfrac12\dfrac{\rho gL^2}{Y}$}{$\dfrac{\rho gL^2}{Y}$}{$\dfrac{1}{4Y}\rho gL^2$}}

\QQ{Water flows through a horizontal pipe of varying cross-section at a rate of $0.314\unit{m^3s^{-1}}$. The velocity of water at a point where the radius of the pipe is $10\unit{cm}$ is}
{\opt{$0.1\unit{m\,s^{-1}}$}{$1\unit{m\,s^{-1}}$}{$10\unit{m\,s^{-1}}$}{$100\unit{m\,s^{-1}}$}}

\QQ{A solid cube of mass $m$ at a temperature $\theta_0$ is heated at a constant rate. It becomes liquid at temperature $\theta_1$ and vapour at temperature $\theta_2$. Let $s_1$ and $s_2$ be specific heats in its solid and liquid states respectively. If $L_f$ and $L_v$ are latent heats of fusion and vaporisation respectively, then the minimum heat energy supplied to the cube until it vaporises is}
{\opt{$ms_1(\theta_1-\theta_0)+ms_2(\theta_2-\theta_1)$}{$mL_f+ms_2(\theta_2-\theta_1)+mL_v$}{$ms_1(\theta_1-\theta_0)+mL_f+ms_2(\theta_2-\theta_1)+mL_v$}{$ms_1(\theta_1-\theta_0)+mL_f+ms_2(\theta_2-\theta_0)+mL_v$}}

\QQ{One mole of an ideal monoatomic gas is taken round the cyclic process MNOM. The work done by the gas is
\fig{\begin{tikzpicture}[scale=0.7]
\draw[->] (0,0) -- (3.8,0) node[right,lbl]{$V$};
\draw[->] (0,0) -- (0,3.4) node[above,lbl]{$P$};
\coordinate (M) at (1,3); \coordinate (N) at (3,1); \coordinate (O) at (1,1);
\draw[densely dashed] (0,3) node[left,lbl]{$3P_0$} -- (M); \draw[densely dashed] (0,1) node[left,lbl]{$P_0$} -- (O);
\draw[densely dashed] (O) -- (1,0) node[below,lbl]{$V_0$}; \draw[densely dashed] (N) -- (3,0) node[below,lbl]{$3V_0$};
\draw[thick,postaction={decorate},decoration={markings,mark=at position 0.55 with {\arrow{Stealth}}}] (M) -- (N);
\draw[thick,postaction={decorate},decoration={markings,mark=at position 0.55 with {\arrow{Stealth}}}] (N) -- (O);
\draw[thick,postaction={decorate},decoration={markings,mark=at position 0.55 with {\arrow{Stealth}}}] (O) -- (M);
\node[lbl,above] at (M) {M}; \node[lbl,right] at (N) {N}; \node[lbl,below left] at (O) {O};
\end{tikzpicture}}}
{\opt{$4.5\,P_0V_0$}{$4\,P_0V_0$}{$9\,P_0V_0$}{$2\,P_0V_0$}}

\QQ{The ratio of molar specific heats of oxygen is}
{\opt{1.4}{1.67}{1.33}{1.28}}

\QQ{For a particle executing simple harmonic motion (SHM), at its mean position}
{\opt{velocity is zero and acceleration is maximum}{velocity is maximum and acceleration is zero}{both velocity and acceleration are maximum}{both velocity and acceleration are zero}}

\QQ{A motor-cyclist moving towards a huge cliff with a speed of $18\unit{km\,h^{-1}}$, blows a horn of source frequency $325\unit{Hz}$. If the speed of the sound in air is $330\unit{m\,s^{-1}}$, the number of beats heard by him is}
{\opt{5}{4}{10}{7}}

\QQ{A body has a charge of $-3.2\unit{\mu C}$. The number of excess electrons it has is}
{\opt{$5.12\times10^{25}$}{$5\times10^{12}$}{$2\times10^{13}$}{$5.12\times10^{13}$}}

\QQ{A point charge A of $+10\unit{\mu C}$ and another point charge B of $+20\unit{\mu C}$ are kept $1\unit{m}$ apart in free space. The electrostatic force on A due to B is $\vec F_1$ and the electrostatic force on B due to A is $\vec F_2$. Then}
{\opt{$\vec F_1 = -2\vec F_2$}{$\vec F_1 = -\vec F_2$}{$2\vec F_1 = -\vec F_2$}{$\vec F_1 = \vec F_2$}}

\QQ{A uniform electric field $E = 3\times10^5\unit{N\,C^{-1}}$ is acting along the positive Y-axis. The electric flux through a rectangle of area $10\unit{cm}\times30\unit{cm}$ whose plane is parallel to the Z-X plane is}
{\opt{$12\times10^3\unit{V\,m}$}{$9\times10^3\unit{V\,m}$}{$15\times10^3\unit{V\,m}$}{$18\times10^3\unit{V\,m}$}}

\QQ{The total electric flux through a closed spherical surface of radius `$r$' enclosing an electric dipole of dipole moment $2aq$ is (Given $\varepsilon_0$ = permittivity of free space)}
{\opt{zero}{$\dfrac{q}{\varepsilon_0}$}{$\dfrac{2q}{\varepsilon_0}$}{$\dfrac{8\pi r^2q}{\varepsilon_0}$}}

\QQ{Under electrostatic condition of a charged conductor, which among the following statements is true?}
{\opt{The electric field on the surface of a charged conductor is $\dfrac{\sigma}{2\varepsilon_0}$, where $\sigma$ is the surface charge density.}{The electric potential inside a charged conductor is always zero.}{Any excess charge resides on the surface of the conductor.}{The net electric field is tangential to the surface of the conductor.}}

\QQ{A cube of side $1\unit{cm}$ contains 100 molecules each having an induced dipole moment of $0.2\times10^{-6}\unit{C\,m}$ in an external electric field of $4\unit{N\,C^{-1}}$. The electric susceptibility of the material is \rule{1.2cm}{0.4pt} $\mathrm{C^2N^{-1}m^{-2}}$.}
{\opt{50}{5}{0.5}{0.05}}

\QQ{A capacitor of capacitance $5\unit{\mu F}$ is charged by a battery of emf $10\unit{V}$. At an instant of time, the potential difference across the capacitor is $4\unit{V}$ and the time rate of change of potential difference across the capacitor is $0.6\unit{V\,s^{-1}}$. Then the time rate at which energy is stored in the capacitor at that instant is}
{\opt{$12\unit{\mu W}$}{$3\unit{\mu W}$}{zero}{$30\unit{\mu W}$}}

\QQ{$\vec E$ is the electric field inside a conductor whose material has conductivity $\sigma$ and resistivity $\rho$. The current density inside the conductor is $\vec j$. The correct form of Ohm's law is}
{\opt{$\vec E = \sigma\vec j$}{$\vec j = \rho\vec E$}{$\vec E = \rho\vec j$}{$\vec E\cdot\vec j = \rho$}}

\QQ{In the circuit shown, the end A is at potential $V_0$ and end B is grounded. The electric current $I$ indicated in the circuit is
\fig{\begin{circuitikz}[scale=0.5,transform shape,font=\Large]
\draw (0,1) node[circ]{} node[above]{$+V_0$} node[below=4pt]{A} -- (0.4,1) to[R,l=$R$] (2.2,1) to[R,l=$R$] (4.0,1) to[R,l=$R$] (5.8,1) to[short,i=$I$] (7.2,1);
\draw (0.4,1) -- (0.4,-0.2) to[R,l_=$R$] (4.0,-0.2) -- (4.0,1);
\draw (2.2,1) -- (2.2,2.2) to[R,l=$R$] (5.8,2.2) -- (5.8,1);
\draw (7.2,1) to[R,l=$2R$] (9.0,1) to[R,l=$2R$] (10.8,1) to[R,l=$2R$] (12.6,1) -- (13.2,1) node[circ]{} node[above]{B} -- (13.2,0) node[ground]{};
\draw (7.2,1) -- (7.2,-0.2) to[R,l_=$2R$] (10.8,-0.2) -- (10.8,1);
\draw (9.0,1) -- (9.0,2.2) to[R,l=$2R$] (12.6,2.2) -- (12.6,1);
\foreach \x in {0.4,2.2,4.0,5.8,7.2,9.0,10.8,12.6} \fill (\x,1) circle (2.5pt);
\end{circuitikz}}}
{\opt{$\dfrac{V_0}{R}$}{$\dfrac{2V_0}{R}$}{$\dfrac{3V_0}{R}$}{$\dfrac{V_0}{3R}$}}

\QQ{The electric current flowing through a given conductor varies with time as shown in the graph below. The number of free electrons which flow through a given cross-section of the conductor in the time interval $0\le t\le20\unit{s}$ is
\fig{\begin{tikzpicture}[xscale=0.11,yscale=0.0055]
\foreach \x in {5,10,...,30} \draw[densely dashed,gray] (\x,0) -- (\x,500);
\foreach \y in {100,200,...,500} \draw[densely dashed,gray] (0,\y) -- (30,\y);
\draw[->] (0,0) -- (33,0) node[right,lbl]{$t$ (s)};
\draw[->] (0,0) -- (0,560) node[above,lbl]{$I$ (in mA)};
\foreach \x in {5,10,15,20,25} \node[lbl,below] at (\x,0) {\x};
\node[lbl,left] at (0,200) {200}; \node[lbl,left] at (0,400) {400};
\draw[line width=0.9pt] (0,100) -- (10,300) -- (30,300);
\fill (0,100) circle (25pt and 500pt);
\end{tikzpicture}}}
{\opt{$3.125\times10^{19}$}{$1.6\times10^{19}$}{$6.25\times10^{18}$}{$1.625\times10^{18}$}}

\QQ{The $I$--$V$ graph for a conductor at two different temperatures $100^\circ$C and $400^\circ$C is as shown in the figure. The temperature coefficient of resistance of the conductor is about (in per degree Celsius)
\fig{\begin{tikzpicture}[scale=0.8]
\draw[->] (0,0) -- (2.8,0) node[right,lbl]{$V$};
\draw[->] (0,0) -- (0,2.2) node[above,lbl]{$I$};
\node[lbl,below left] at (0,0) {O};
\draw (0,0) -- (45:2.4); \draw (0,0) -- (30:2.6);
\draw (0.75,0) arc (0:30:0.75); \node[lbl] at (1.0,0.22) {$30^\circ$};
\draw (30:1.35) arc (30:45:1.35); \node[lbl] at (37:1.65) {$45^\circ$};
\end{tikzpicture}}}
{\opt{$3\times10^{-3}$}{$6\times10^{-3}$}{$9\times10^{-3}$}{$12\times10^{-3}$}}

\QQ{An electric bulb of $60\unit{W}$, $120\unit{V}$ is to be connected to $220\unit{V}$ source. What resistance should be connected in series with the bulb, so that the bulb glows properly?}
{\opt{$50\unit{\Omega}$}{$100\unit{\Omega}$}{$200\unit{\Omega}$}{$288\unit{\Omega}$}}

\QQ{In an experiment to determine the temperature coefficient of resistance of a conductor, a coil of wire X is immersed in a liquid. It is heated by an external agent. A meter bridge set up is used to determine resistance of the coil X at different temperatures. The balancing points measured at temperatures $t_1 = 0^\circ$C and $t_2 = 100^\circ$C are $50\unit{cm}$ and $60\unit{cm}$ respectively. If the standard resistance taken out is $S = 4\unit{\Omega}$ in both trials, the temperature coefficient of the coil is
\fig{\begin{circuitikz}[scale=0.55,transform shape,font=\Large]
\draw[line width=1.2pt] (0,1.6) -- (0,2.4) -- (2.5,2.4) (3.5,2.4) -- (6.3,2.4) (7.3,2.4) -- (9.6,2.4) -- (9.6,1.6);
\draw (0,1.6) -- (9.6,1.6);
\draw (1.0,2.4) -- (1.0,4.4) -- (1.5,4.4); \draw (3.0,4.4) -- (3.5,4.4) -- (3.5,2.4);
\draw (1.4,3.3) -- (1.4,4.6) (3.1,3.3) -- (3.1,4.6) (1.4,3.3) -- (3.1,3.3);
\fill[pattern=dots] (1.45,3.35) rectangle (3.05,4.2);
\draw (1.6,4.4) -- (1.9,4.4) to[L] (2.6,4.4) -- (2.9,4.4); \node at (2.25,4.85) {X};
\draw (2.25,3.3) -- (2.25,2.9) node[ground]{};
\draw (6.3,2.4) -- (6.3,3.8) to[R,l={$S=4\,\Omega$}] (8.3,3.8) -- (8.3,2.4) (7.3,2.4) -- (8.3,2.4);
\draw (4.9,2.4) to[rmeter,t=G] (5.6,1.6);
\draw[<->] (0,1.2) -- (5.6,1.2) node[midway,below]{$l$};
\draw (0,1.6) -- (0,0) to[battery1] (5,0) to[switch] (9.6,0) -- (9.6,1.6);
\end{circuitikz}}}
{\opt{$0.05^\circ\mathrm{C^{-1}}$}{$0.02^\circ\mathrm{C^{-1}}$}{$0.005^\circ\mathrm{C^{-1}}$}{$2.0^\circ\mathrm{C^{-1}}$}}

\QQ{A moving electron produces}
{\opt{only electric field}{both electric and magnetic field}{only magnetic field}{neither electric nor magnetic field}}

\QQ{A coil having 9 turns carrying a current produces magnetic field $B_1$ at the centre. Now the coil is rewounded into 3 turns carrying same current. Then the magnetic field at the centre $B_2 =$ \rule{1.2cm}{0.4pt}.}
{\opt{$\dfrac{B_1}{9}$}{$9B_1$}{$3B_1$}{$\dfrac{B_1}{3}$}}

\QQ{A particle of specific charge $\dfrac qm = \pi\unit{C\,kg^{-1}}$ is projected from the origin towards positive $x$-axis with the velocity $10\unit{m\,s^{-1}}$ in a uniform magnetic field $\vec B = -2\hat k\unit{T}$. The velocity $\vec v$ of particle after time $t = \dfrac1{12}\unit{s}$ will be (in $\mathrm{m\,s^{-1}}$)}
{\opt{$5(\hat i+\hat j)$}{$5(\hat i+\sqrt3\,\hat j)$}{$5(\sqrt3\,\hat i-\hat j)$}{$5(\sqrt3\,\hat i+\hat j)$}}

\QQ{The magnetic field at the centre of a circular coil of radius $R$ carrying current $I$ is 64 times the magnetic field at a distance $x$ on its axis from the centre of the coil. Then the value of $x$ is}
{\opt{$\dfrac R4\sqrt{15}$}{$R\sqrt3$}{$\dfrac R4$}{$R\sqrt{15}$}}

\QQ{Magnetic hysterisis is exhibited by \rule{1.2cm}{0.4pt} magnetic materials.}
{\opt{only para}{only dia}{only ferro}{both para and ferro}}

\QQ{Magnetic susceptibility of Mg at $300\unit{K}$ is $1.2\times10^{-5}$. What is its susceptibility at $200\unit{K}$?}
{\opt{$18\times10^{-5}$}{$180\times10^{-5}$}{$1.8\times10^{-5}$}{$0.18\times10^{-5}$}}

\QQ{A uniform magnetic field of strength $B = 2\unit{mT}$ exists vertically downwards. These magnetic field lines pass through a closed surface as shown in the figure. The closed surface consists of a hemisphere $S_1$, a right circular cone $S_2$ and a circular surface $S_3$. The magnetic flux through $S_1$ and $S_2$ are respectively
\fig{\begin{tikzpicture}[scale=0.8]
\draw (-1.3,0) arc (180:0:1.3);
\draw (-1.3,0) arc (180:360:1.3 and 0.25); \draw[dashed] (1.3,0) arc (0:180:1.3 and 0.25);
\draw (-1.3,0) -- (0,-2.0) -- (1.3,0);
\fill (0,0) circle (1.2pt); \draw[dashed,->] (0,0) -- (1.3,0);
\node[lbl,right] at (1.4,0) {$R = \dfrac{10}{\sqrt\pi}$};
\foreach \x in {-0.55,0.55} \draw[->] (\x,2.0) -- (\x,-2.3);
\node[lbl] at (1.15,1.15) {$S_1$}; \node[lbl] at (-0.35,0.45) {$S_3$}; \node[lbl] at (1.0,-0.95) {$S_2$};
\node[lbl] at (0.85,-2.1) {$\vec B$};
\end{tikzpicture}}}
{\opt{$\Phi_{S_1} = -20\unit{\mu Wb}$, $\Phi_{S_2} = +20\unit{\mu Wb}$}{$\Phi_{S_1} = +20\unit{\mu Wb}$, $\Phi_{S_2} = -20\unit{\mu Wb}$}{$\Phi_{S_1} = -40\unit{\mu Wb}$, $\Phi_{S_2} = +40\unit{\mu Wb}$}{$\Phi_{S_1} = +40\unit{\mu Wb}$, $\Phi_{S_2} = -40\unit{\mu Wb}$}}

\QQ{In the figure, a conducting ring of certain resistance is falling towards a current carrying straight long conductor. The ring and conductor are in the same plane. Then the
\fig{\begin{tikzpicture}[scale=0.8]
\draw (0.6,1.3) circle (0.5); \node[lbl] at (1.15,1.9) {A};
\draw[->] (-0.2,1.7) -- (-0.2,0.9);
\draw[densely dashed] (-1.6,0) -- (2.6,0);
\fill[gray!70] (-0.9,-0.07) rectangle (1.7,0.07); \node[lbl,above] at (1.3,0.1) {B};
\draw[->] (0.8,-0.35) -- (1.6,-0.35) node[midway,below,lbl]{$I$};
\end{tikzpicture}}}
{\opt{induced electric current is zero}{induced electric current is anticlockwise}{induced electric current is clockwise}{ring will come to rest}}

\QQ{An induced current of $2\unit{A}$ flows through a coil. The resistance of the coil is $10\unit{\Omega}$. What is the change in magnetic flux associated with the coil in $1\unit{ms}$?}
{\opt{$0.2\times10^{-2}\unit{Wb}$}{$2\times10^{-2}\unit{Wb}$}{$22\times10^{-2}\unit{Wb}$}{$0.22\times10^{-2}\unit{Wb}$}}

\QQ{A square loop of side length `$a$' is moving away from an infinitely long current carrying conductor at a constant speed `$v$' as shown. Let `$x$' be the instantaneous distance between the long conductor and side AB. The mutual inductance (M) of the square loop -- long conductor pair changes with time ($t$) according to which of the following graphs?
\fig{\begin{tikzpicture}[scale=0.75]
\fill[gray!70] (-0.08,-0.6) rectangle (0.08,2.0); \draw[densely dashed] (0,-1.2) -- (0,-0.6) (0,2.0) -- (0,2.6);
\draw[->] (0.3,1.6) -- (0.3,2.3) node[right,lbl]{$I$};
\draw (1.0,0) rectangle (2.2,1.2);
\node[lbl,left] at (1.0,1.2) {A}; \node[lbl,left] at (1.0,0) {B}; \node[lbl,right] at (2.2,0) {C}; \node[lbl,right] at (2.2,1.2) {D};
\node[lbl,above] at (1.6,1.2) {$a$}; \node[lbl,below] at (1.6,0) {$a$}; \node[lbl,right] at (1.0,0.6) {$a$}; \node[lbl,left] at (2.2,0.6) {$a$};
\draw[->] (1.9,0.9) -- (2.7,0.9) node[right,lbl]{$v$};
\draw[<->] (0.08,-0.2) -- (1.0,-0.2) node[midway,below,lbl]{$x$};
\end{tikzpicture}}}
{\optfig{\gaxes{$t$}{$M$}{\draw (0.2,1.6) .. controls (0.5,0.5) and (1.0,0.25) .. (1.7,0.15);}}
{\gaxes{$t$}{$M$}{\draw (0,0) -- (1.8,1.1);}}
{\gaxes{$t$}{$M$}{\draw (0,1.1) -- (1.7,1.1);}}
{\gaxes{$t$}{$M$}{\draw (0,1.6) -- (1.2,0);}}}

\QQ{Which of the following combinations should be selected for better tuning of an LCR circuit used for communication?}
{\opt{$R = 20\unit{\Omega}$, $L = 1.5\unit{H}$, $C = 35\unit{\mu F}$}{$R = 25\unit{\Omega}$, $L = 2.5\unit{H}$, $C = 45\unit{\mu F}$}{$R = 25\unit{\Omega}$, $L = 1.5\unit{H}$, $C = 45\unit{\mu F}$}{$R = 15\unit{\Omega}$, $L = 3.5\unit{H}$, $C = 30\unit{\mu F}$}}

\QQ{In an LCR series circuit, the value of only capacitance $C$ is varied. The resulting variation of resonance frequency ($f_0$) as a function of $C$ can be represented as}
{\optfig{\gaxes{$C$}{$f_0$}{\draw (0.1,0.1) .. controls (0.6,0.1) and (0.6,1.5) .. (0.9,1.5) .. controls (1.2,1.5) and (1.2,0.1) .. (1.7,0.1);}}
{\gaxes{$C$}{$f_0$}{\draw (0,0) -- (1.6,1.6);}}
{\gaxes{$C$}{$f_0$}{\draw[domain=0.12:1.9,samples=40] plot (\x,{0.5/sqrt(\x)-0.15});}}
{\gaxes{$C$}{$f_0$}{\draw[domain=0:1.9,samples=40] plot (\x,{1.2*sqrt(\x)});}}}

\QQ{The figure shows variation of $R$, $X_L$ and $X_C$ with frequency `$f$' in a series LCR circuit. Then for what frequency point is the circuit capacitive?
\fig{\begin{tikzpicture}[scale=0.8]
\draw[->] (0,0) -- (4.3,0) node[right,lbl]{$f$};
\draw[->] (0,0) -- (0,2.8);
\draw (0,0.9) -- (3.6,0.9) node[above,lbl]{$R$};
\draw (0,0) -- (2.3,2.4) node[above,lbl]{$X_L$};
\draw[domain=0.3:2.8,samples=40] plot (\x,{0.75/\x+0.15}) ; \node[lbl,above] at (0.32,2.5) {$X_C$};
\draw[densely dashed] (0.86,0.9) -- (0.86,0) node[below,lbl]{A};
\draw[densely dashed] (1.12,1.17) -- (1.12,0) node[below,lbl]{B};
\draw[densely dashed] (1.25,0.9) -- (1.25,0) node[below,lbl]{C};
\draw[densely dashed] (2.0,0.53) -- (2.0,0) node[below,lbl]{D};
\end{tikzpicture}}}
{\opt{B}{D}{A}{C}}

\QQ{Electromagnetic waves are incident normally on a perfectly reflecting surface having surface area $A$. If $I$ is the intensity of the incident electromagnetic radiation and $c$ is the speed of light in vacuum, the force exerted by the electromagnetic wave on the reflecting surface is}
{\opt{$\dfrac{2IA}{c}$}{$\dfrac{IA}{c}$}{$\dfrac{IA}{2c}$}{$\dfrac{I}{2Ac}$}}

\QQ{The final image formed by an astronomical telescope is}
{\opt{real, erect and diminished}{virtual, inverted and diminished}{real, inverted and magnified}{virtual, inverted and magnified}}

\QQ{If the angle of minimum deviation is equal to angle of a prism for an equilateral prism, then the speed of light inside the prism is \rule{1.2cm}{0.4pt}.}
{\opt{$3\times10^8\unit{m\,s^{-1}}$}{$2\sqrt3\times10^8\unit{m\,s^{-1}}$}{$\sqrt3\times10^8\unit{m\,s^{-1}}$}{$\dfrac{\sqrt3}{2}\times10^8\unit{m\,s^{-1}}$}}

\QQ{A luminous point object O is placed at a distance $2R$ from the spherical boundary separating two transparent media of refractive indices $n_1$ and $n_2$ as shown, where $R$ is the radius of curvature of the spherical surface. If $n_1 = \frac43$, $n_2 = \frac32$ and $R = 10\unit{cm}$, the image is obtained at a distance from P equal to
\fig{\begin{tikzpicture}[scale=0.75]
\draw[decorate,decoration={snake,amplitude=1pt,segment length=6pt}] (0,-1.3) -- (0,1.3);
\draw[decorate,decoration={snake,amplitude=1pt,segment length=6pt}] (5.2,-1.3) -- (5.2,1.3);
\draw (0,1.3) -- (5.2,1.3) (0,-1.3) -- (5.2,-1.3);
\draw (2.8,1.3) arc (110:250:1.38);
\draw[densely dashed] (0,0) -- (5.2,0);
\fill (0.3,0) circle (1.6pt) node[above,lbl]{O}; \fill (3.8,0) circle (1.6pt) node[above,lbl]{C};
\node[lbl,below right] at (2.33,0) {P};
\draw[<->] (2.33,0.55) -- (3.8,0.55) node[midway,above,lbl]{$R$};
\draw[<->] (0.3,-0.85) -- (2.33,-0.85) node[midway,below,lbl]{$2R$};
\node[lbl] at (1.0,0.85) {$n_1 = \frac43$}; \node[lbl] at (4.2,0.95) {$n_2 = \frac32$};
\end{tikzpicture}}}
{\opt{$30\unit{cm}$ in the rarer medium}{$30\unit{cm}$ in the denser medium}{$18\unit{cm}$ in the rarer medium}{$18\unit{cm}$ in the denser medium}}

\end{multicols}
